% Generated from ../chapters/26-the-organism-as-an-open-question.md
\chapter{The Organism as an Open Question}

The body is neither a passive machine nor an unlimited miracle.

It is a historically constrained, materially vulnerable, self-organizing process. It contains remarkable repair, learning, and adaptation, together with hard limits, trade-offs, and tissue-specific failures.

The framework developed here begins with the river: restore dynamics as well as parts. It adds downward causation: beliefs, projects, attention, and behaviour become biological through ordinary cascades. It proposes the healthy attractor: development constrains the forms toward which coordinated function can converge. It treats relaxation as the release of degrees of freedom and progressive challenge as a way to expand the system's repertoire. It interprets aging as a reinforcing loss of resilience and asks whether improvement can also become self-reinforcing.

The strongest claim---endogenous maintenance escape velocity---remains unproven. It may be wrong, partially right, or applicable only to selected tissues and timescales.

But the practical framework does not require certainty.

Live as though adaptive capacity is precious.

Challenge it without crushing it.

Release control without abandoning structure.

Alternate building with cleanup and recovery.

Measure durable function rather than dramatic sensation.

Use external repair without surrendering the organism's role in healing.

Most importantly, remain open to evidence.

The self-organizing body is not a conclusion. It is a research programme---and a way of living inside the question.

\begin{center}\rule{0.5\linewidth}{0.5pt}\end{center}
